\itemhead{3.1 思考 · 未编号题组（教材p.65）}{XJ-B1-C03-TH-P065-Y0594}
\label{XJ-B1-C03-TH-P065-Y0594}
{\footnotesize 湘教版教材p.65 / PDF第72页 / 3.1 思考 / 未编号题组}
\field{原题}{在初中，函数概念是怎样定义的？与函数概念有关的术语有哪些？}
\field{参考答案}{初中从变量依赖关系定义函数：在一个变化过程中，对自变量每个确定的值，因变量都有唯一确定的值与之对应。相关术语为自变量、因变量、函数值、定义域、值域。}
{\small \field{关联知识}{F01 函数本质、构成要素与相同函数（教材帮p.77、78 / PDF86、87）}
\field{方法与易错点}{逐一核对每个允许输入是否有且仅有一个合乎目标集合的输出。 易错：把值域直接当作目标集合；只看解析式外观；认为常函数不是函数。}
\field{答案依据}{答案册未提供本栏目的独立对应答案；依据原题、教材概念补写参考回应。开放性总结不设唯一数值答案。}}
